\subsection*{Strictly scalar local action learning}

The primary action comparison uses the same conservative transport physics, with initial mass 24 and energy 144. Two noise seeds, 202610199431 and 202610199432, are paired with unknown physical rotations $+90^\circ$ and $-90^\circ$. The fixed periodic substrate is $0.6\sin(x)+0.4\sin(y)$ in the environment's spatial convention. Its coordinates and values are not policy inputs. A site with $\rho\geq10^{-4}$ can receive a completed payment; positive transfer is limited by local energy capacity and negative transfer by available energy. The policy observes
\begin{equation}
y_i=\frac{G_i-L_i}{0.005\rho_i}
\end{equation}
only where a nonzero transfer occurred. Normalization does not divide by field amplitude. Missing payments remain missing observations, not zero or adverse outcomes. The learning objective uses these environmental receipts; movement expenditure is reported separately and is not subtracted again to define the learner's scalar signal.

An internal eight-update oscillator controls exploration, without exposing the environmental schedule. At the start of each action block, the command is $\boldsymbol u=0.8\tanh(\boldsymbol m+0.6\boldsymbol\epsilon)$, with two independent standard Gaussian innovations. The command and score $\boldsymbol s=\boldsymbol\epsilon/0.6$, the starting scalar receipt and its validity flag travel with the material. An additional independent random heading field is generated on the same budget for all policies. The physical operator $K$ carries every extensive internal quantity using the same actual transported density. Its entries, donor displacement, velocity, spatial reward differences and the hidden actuator map are never passed to \texttt{choose} or \texttt{observe}. A transported indicator records whether actual motion was paid; this indicator supplies no motion direction.

At a completed block, a response update requires a current paid outcome, a valid transported starting observation and a completed-motion receipt. Using the transported score and starting observation, the fixed-rate response estimator is
\begin{equation}
\boldsymbol g\leftarrow(1-\beta)\boldsymbol g+
\beta\boldsymbol s\,(y_{\mathrm{end}}-y_{\mathrm{start}}),
\qquad \beta=1/8,\qquad \boldsymbol m=\operatorname{unit}(\boldsymbol g).
\end{equation}
The estimate starts at zero; zero norm gives zero mean. The rate was fixed before execution and not searched. The historical code label \texttt{spsa} denotes this single-trajectory perturbation-response rule. There are no paired simultaneous counterfactual evaluations, and forward material mixing prevents an interpretation as an unbiased whole-organism policy gradient. The rule estimates a reward-related action response, not an identified $2\times2$ mechanical mapping.

The accumulated policy updates a constant two-component mean with rate 2. Its context variant additionally reads its own current scalar and the previous completed block's scalar change; features are normalized and coefficients clipped to $[-2,2]$. Run-and-tumble retains its direction after an improving block and chooses a random heading after deterioration. The random policy writes no acquired direction. Response normalization, fixed-policy unit headings, command bound, noise and payment rules remain explicit: equal tariffs and exploration opportunities do not imply equal realized forces or costs. State storage and learning arithmetic have no additional physical tariff.

Life-zero trajectories comprise 320 acquisition steps, 320 steps with the reward meaning reversed and 128 steps with the same meaning but no new payments, reward observations or parameter writes. A phase passes only if its final-64-step mean field score exceeds 0.25, favourable mass fraction is at least 0.6, total mass is at least 12 and the largest connected fraction is at least 0.75. The full recorded stage and its failed criterion are retained even when a later phase passes. The two response-estimation lives are actual new runs under the frozen seeds and schedule; earlier four-policy data are reused read-only. 

Content probes restore the complete actual step-640 parents, preserve density, energy, other registrations, physical coupling, RNG and clock, then clear $\boldsymbol g$ or rotate it by $90^\circ$. All subsequent writes and outcome payments are disabled for 128 steps. Recorded retained tails supply the reference. The execution-change assay restores the complete step-320 parents; at step 321 the physical $\pm90^\circ$ map becomes the opposite $45^\circ$ map while reward meaning stays fixed. Each branch receives 320 paid steps and a 128-step aftereffect. The fixed run-and-tumble control is initialized with $\operatorname{unit}(\boldsymbol g)$ from that parent. Raw phase labels retain the recorder's word \texttt{reversal}, but this assay changes execution, not reward meaning. Frozen estimation stops only $\boldsymbol g$ updates: movement, lawful scalar feedback and transport continue. Source rows distinguish these interventions rather than pooling them under a common reversal success count.

\subsection*{Continuing-body cue assay}

The cue assay uses a periodic $32\times32$ lattice with density $\rho$, intracellular energy $E$, two cue values $Q_i$, a two-channel eligibility field $z_i$ and the most recently registered cue $\ell_i$. Initial total mass and energy are 24 and 72; cue values start at $0.2$. Two density-perturbation seeds, 202610197401 and 202610197402, are crossed with both initial cue--feedback orders. These four constructed conditions are evaluated with reactor, recent-feedback and eligibility policies under paired and unpaired exposures. There is no outer optimization in this batch.

The inherited material-transport model uses local four-neighbour movement, cohesion and pressure. Its fixed coefficients include time step $0.1$, cohesion 3, pressure 1.5, spatial scale 1.4 and cue gain 18. Local directed rates use a positive, bounded response to neighbour potential differences and to cue differences weighted by local values. The actual transported density and energy determine subsequent state; the external environment does not translate or reset the body between lessons. Movement and basal prices are $0.002$ and $0.012$ in the model's native units, with donor-energy limitations and capacity $4\rho$. The source implementation and complete coefficient file are retained rather than substituting a point-agent approximation for these dynamics.

At each actual transport, extensive amounts $\rho Q$, $\rho z$ and $\rho\ell$ travel with the recorded material operator $K$. Their intensive values are recovered using the corresponding transported mass, with zero-density entries masked. This registration is different from simply copying a value at its original lattice coordinate. Tissue may mix at a destination, so the resulting intensive average does not preserve distinct identities of all microscopic sources. No regenerative births occur in the cue assay.

Each step first executes the physical body update with its built-in plasticity disabled. The experiment then applies the prescribed local energy feedback to the resulting tissue. If $f$ is the signed external feedback field, the requested magnitude is $r=0.08\rho|f|$. Positive feedback deposits $g=\min(r,\max(4\rho-E,0))$; negative feedback removes $d=\min(r,E)$. Only the term appropriate to the sign is applied. The realized teaching signal is
\begin{equation}
y=\frac{g-d}{r},
\end{equation}
with zero where the denominator is absent. It therefore reports actual local delivered gain or loss, not the hidden environmental cue mapping. Normalization also means that very small nonzero feedback in a Gaussian tail can remain recognizable; this assay does not establish difficult credit assignment under a finite detection threshold.

All policies maintain the same sensory registrations. Where total current cue exceeds $0.05$, the normalized current cue replaces $\ell$. The eligibility field follows
\begin{equation}
z_i\leftarrow e^{-1/16}z_i+(1-e^{-1/16})c_i.
\end{equation}
Recent feedback uses $\phi=\ell$ and $\alpha=1$; eligibility uses $\phi=z/\sum_i z_i$ and $\alpha=0.5$. At sites with nonzero realized feedback and positive tissue, the update is
\begin{align}
\delta &=y-\sum_i Q_i\phi_i,\\
Q_i&\leftarrow\operatorname{clip}\left(Q_i+
\frac{\alpha\phi_i\delta}{\max(\sum_j\phi_j^2,10^{-12})},-1,1\right).
\end{align}
The reactor suppresses these writes while maintaining cue registrations and transport. There is no policy-gradient estimate\cite{williams,baxter}, no inferred outcome for an unobserved action and no learner input for phase, time, seed, coordinates or reversal identity. The supplied movement rule reads local values and local cue gradients. Extra learning arithmetic has no separately assigned physical price.


\subsection*{Exposure, reversal, probing and value erasure}

Each life lasts 960 updates. An initial 64-step probe is followed by two acquisition epochs. Each epoch contains a 48-step lesson for each cue and a 64-step unreinforced probe. A 64-step neutral interval and 64-step retention probe follow. The environmental mapping then reverses without informing the learner; two further lesson--probe epochs, another neutral interval and another retention probe complete the life. Probe endpoints are 64, 224, 384, 512, 672, 832 and 960.

During a lesson, a Gaussian cue of scale 6 is centred at the body's position at lesson onset and remains fixed there for that lesson. Cue presentation occupies the first 24 steps; four feedback pulses occur at offsets 32, 36, 40 and 44. Paired exposure uses a consistent sign for each cue within the current mapping. Unpaired exposure balances four positive and four negative pulses per cue across the two epochs, using the recorded seed shuffle in both mapping orders. Exposure is broad and imposed, rather than earned through free foraging. The environment uses the current centroid to place a new trial; the learner does not receive that centroid.

Probes place the two cue Gaussians at opposite offsets of four cells with scale 3.5, changing left--right placement across probes. No outcome feedback is delivered. Let $A_i$ be actual density integrated over the radius-2.5 contact region around cue $i$. The primary probe score is
\begin{equation}
p_{\rm correct}=\frac{A_{\rm currently\ positive}}{A_0+A_1}.
\end{equation}
The score is undefined when its denominator lacks support; no missing value is replaced by a successful response. We report endpoint fractions, condition-wise ranges and descriptive condition means. They are distinct from time-averaged probe contact or signed centroid displacement, which remain in the source tables.

The reversal occurs at the 512 boundary. The first new completed feedback is at 545, and the first post-reversal behavioural probe spans 609--672. By that probe the body has received four feedback pulses for each cue. Recovery is therefore established by the first measured post-exposure probe, not after one feedback event. Choice during forced exposure is not reconstructed as an independent error rate.

Each 960-step terminal state supplies two 64-step unreinforced continuations. Both disable value writes. The retained branch keeps $Q$, and the erased branch sets $Q_i=0.2$ while preserving density, energy, eligibility and last-cue records exactly at branching. Transport and sensing continue. These paired forks measure expression under the remaining bodily state; they do not erase the entire acquisition history. Viability requires mass at least 12 and a largest connected component containing at least half the mass, without 32 consecutive failing updates. The final report contains 24 main lives and 48 forks, all alive; pilot trajectories are separate. Its total accounting includes 28,032 actual updates across science and pilots, whereas the 72 scientific paths comprise 26,112 updates. No trajectory was rerun for this manuscript's summary.


\subsection*{Distractor histories and changing temporal sources}

The temporal-credit experiment uses two additional seeds, 202610197403 and 202610197404, with both initial cue-meaning orders. Twenty lives begin under five rules. Each 40-step lesson presents an earlier cue for eight steps, a four-step gap, a recent cue for eight steps and another four-step gap; feedback is delivered at completed updates 25, 29, 33 and 37. Each four-lesson block traverses the full earlier--recent cue cross in randomized order. Local cue onsets update two event registers, which travel with the material through the same actual transport as the body. All arms maintain the same auxiliary fields; their rules differ in how those available records are used. Neither the temporal source nor a hidden cue label is passed to a learner.

Let $r_0,r_1$ be the two-component recent and earlier cue registers, and $b_0,b_1$ their associated coefficient vectors. Recent-cue, fixed-trace and fixed-earlier delta use the normalized update above with learning rate one and the respective feature. The fixed trace retains the exponential time constant of sixteen updates. Joint delta uses the concatenated four-component feature and prediction $\hat y=\sum_{j=0}^1 b_j\cdot r_j$ in a normalized least-mean-squares update. For gated experts, the two pre-update predictions are $\hat y_j=b_j\cdot r_j$, and each expert receives its own normalized unit-rate delta update. At sites with nonzero realized feedback, the earlier-source log weight changes according to
\begin{equation}
g\leftarrow\operatorname{clip}\bigl(0.99g+0.75[(\hat y_0-y)^2-(\hat y_1-y)^2],-8,8\bigr),
\qquad w=(1+e^{-g})^{-1}.
\end{equation}
The motor readout uses $(1-w)b_0+wb_1$ for the gated rule and $b_0+b_1$ for joint delta, clipped to $[-1,1]$. The weight is updated from prediction errors computed before the coefficient updates; it is not supplied with the true lag. These calculations add no native bodily step and carry no separately priced arithmetic or storage cost.

The initial distractor experiment lasts 1216 steps for surviving bodies and includes six fixed acquisition, cue-meaning reversal and retention probes. A complete behavioural pass requires each endpoint contact fraction to be at least $0.65$ and the body to survive. The three connectivity failures stop their lives; unobserved probes and terminal forks are not generated. At clock 1216, the seventeen surviving full states are loaded exactly for a 576-step continuation. Cue meanings remain fixed while the recent event replaces the earlier event as feedback source. Each rule continues its own actual prior body and state. The first changed-source feedback is at continuous clock 1241 and the first new probe spans 1377--1440, after sixteen pulses. Two four-lesson blocks with probes are followed by a neutral interval and retention probe. Endpoint contact fractions use the radius-2.5 regions defined above, not Gaussian-cue intensity weights or time-averaged accuracy.


\subsection*{Spatial writing and matched-state interventions}

The earlier spatial-writing example branches at step 240 from identical acquired body and remaining world. Stopping new writing leaves decay, transport, eligibility and subsequent local movement active. Saved fields use original periodic coordinates and common scales. Density masks specify tissue support, not confidence in an inferred value. The local-state and specific-energy panels are observations of evolving fields; they are not separate interventions on those fields. The full 61-point path is retained along with distinct injury and starvation accounts. This model also has no birth. Full source-specific formulas appear in the accompanying detailed protocols.


\subsection*{Resource-constrained generation}

In the generative models, a signed rule response separates deletion from requested production. Positive production is limited by spatial capacity, a local template and actual precursor P and fuel F. The admission order is retained, so a conditional downstream rejection is not equated with an independent deficit evaluated before upstream restriction. Negative response removes actual tissue and its associated energy. Resource transport, uptake, pumping, deletion and conversion have the costs and sinks recorded by that model. External interventions that coordinate requests or exchange resource packages are labelled explicitly; their extra computation or external work is not silently charged to the organism.


\subsection*{Density-regulated regeneration with energetic consequences}

This host uses actual material transport and the existing paid local conversion constraints, while replacing the kernel-derived generative request with a fixed density rule. With physical step $dt=0.05$, rule death removes $0.005$ times the post-actor body, and the requested birth is $0.015$ times the local $3\times3$ binomial-weighted survivor average, with weights $[1,2,1]\otimes[1,2,1]/16$, times $[1-\rho_{\mathrm{survivor}}/0.8]_+$. Actual admission still obeys local $P/F$ support and energy prices. Old dual-kernel metadata supplies grid size and clock but not the generative rate. A uniform external bath restores local $P$ and $F$ to $0.2$ before each native step; actual additions are logged. The apparatus and central computation are unpriced, so this is not a closed metabolism.

Two initial geometry recipes, compact and elongated, each have both initial outcome meanings, with shared within-condition physical states and exploration tapes for fast and fixed-prior controls. Founder energy is $1.8M$ in both controls. Two Gaussian cues at $(x,y)=(26,32)$ and $(38,32)$ use widths $5$ and $12$ with an explicit winner-take-local partition, avoiding simultaneous cue aliasing. A central five-action readout and fixed motor gain $0.45$ are installed. Bounded current-cue learning uses $\alpha=1$; the prior control never writes acquired values but retains sensing, transport, maintenance and copying. Actual uptake is followed, before generation, by a local energy tax of $0.4$ per physical time per old mass in the relevant region. Only the affordable amount enters heat; feedback uses actual uptake minus this paid tax. Unpaid tax is not fictitious negative feedback, and cue identity does not directly delete tissue.

The main schedule is acquisition through 256, unannounced outcome reversal through 512, then stop-writing expression through 640 with physical outcomes continuing. Four actual fast-256 parents also supply paired frozen-writing reversal continuations to 512; the matching fast futures are reused read-only. Survival requires $M\ge0.5M_0$, a largest connected component (density threshold $0.02$) of at least $0.8$ of body mass, and at most $0.01$ in the two-cell border. The final 64 updates of a behavioural phase require mean safe-contact mass at least 8, safe fraction at least $0.65$, and actual new-tissue fraction at least $0.1$. The contact measure is specific to the exclusive cue regions and differs from both the local-motor field score and the earlier radius-2.5 contact assay. Missing terminal windows have no evaluated fraction or pass outcome in the manuscript tables.

Only the two acquired compact states qualify for the newborn-content probe. From the same complete actual-256 parent, writers stop and meaning stays unchanged for 128 updates. At actual birth, copied pre-write local contents are retained, zeroed under the same tariff, or swapped between cue channels with their norm preserved. Zeroing is not norm matched. Old tissue values remain unchanged at the intervention; newborn personal registrations, seen flags, traces and pending state begin at zero. Holding requests $0.00005$ energy units per unit of old mass per step. Copying requests $0.002$ energy units per unit of actually admitted newborn mass; realized payments are energy-limited. Denied charges remain explicit, and realized later fees can differ as futures diverge. Sensing, physical outcomes and transport continue. Material labels record tissue actually born since life zero and since the probe branch; they are not estimates of one-to-one founder replacement when body mass also grows.






\subsection*{Predictive maintenance and equal-purchase physical controls}

The maintenance series uses the explicitly supplied density-regulated birth/death host described above. Local precursor and fuel are restored to 0.2 before each update, with actual additions logged. Protection transfers real precursor into a reserve, costs 0.3 units of old intracellular energy per unit synthesized, incurs holding cost 0.001 per reserve unit per step and decays by 3\% per step. Total local capacity is 0.08 times current tissue, local production is limited to 0.008 times tissue per step, and the initial common production rate is 0.004 times total mass. Realized production is further limited by precursor and payable old energy. Loss of reserve and tissue is registered in waste and dead-energy accounts.

Two neutral warning channels are shown for eight steps, followed by two silent steps and four steps of actual mechanical contact. In the initial jittered protocol, contact onsets are separated by 22--27 steps. Spatial halves receive weak strain 0.002 or strong strain drawn from 0.012/0.024; cue identity predicts which side receives stronger contact. The learner reads current warning and completed incident mechanical strain, including strain absorbed by protection. This is supplied mechanical sensing, not inference of a counterfactual unprotected injury. Tissue demand is $\rho\,\mathrm{strain}$, absorbed amount is the lesser of reserve and demand, and residual demand removes actual tissue and its energy. Damage precedes learning and new purchasing. No future strain map, relation-change flag or absolute environment clock enters the learner.

Local cue registrations and conditional estimates move with actual material transport. Standard delta learning updates two conditional means only at actual contact. Its four charged scalars include the registrations; the risk variant additionally stores second moments and pays for six. Newborn registrations start at zero, while pre-write conditional content is copied from surviving donors subject to the actual charge. Earlier incomplete runs and unmatched-purchase controls are retained in the source records; only the completed, actually matched assays support the equal-purchase comparisons.

The subsequent equal-purchase assay fixes each per-step quota at 0.0015 times registered parent mass and distributes it using $\rho(0.002+\text{conditional estimate})$ subject to the same capacity and payment constraints. Correct, pooled and swapped readouts retain the same stored quantities and tariffs. All six branches actually realize the shared quota; their three 96-step futures per acquired parent total 576 updates. Equal synthesis is not equal total expense, because holding, movement, copying and births respond to each physical future.

The strong-physiology assay restores two actual delta step-128 parents and applies six policies for 192 steps each. The warning relationship changes after the first 96 steps. This follow-up uses fixed 24-step onset spacing, with four contacts before and four after the change (warning starts 132, 156, 180, 204, 228, 252, 276 and 300). Continued delta, frozen content, each fixed mapping, immediate response to completed strain and uniform allocation share the same quota allocator and four charged scalars. Both fixed mappings are evaluated unchanged in both worlds. Their redundant predictor capacity is an explicit charge to simple controls, precluding claims of economic dominance over a minimal implementation. A passive cohort labels all tissue present at branching and is transported and removed by the actual physical operators; newborns never enter this registered old cohort. Stage deaths are normalized by registered parent mass, not by surviving tissue alone.

The timing assay uses two actual step-192 parents from that strong-control series. Frozen learned and instantaneous readouts each receive a 96-step future under distributed four-step contact or a single concentrated contact with the same per-cell integrated external strain. Warning starts are 196, 220, 244 and 268: four contacts per 96-step future at fixed 24-step spacing. Onset and warning are matched, and prediction remains frozen. The impulse applies four times the strain only at warning-start plus ten, replacing the four-step contact block. The environment changes the time profile; the learner senses only current stimuli. Equal external strain does not imply equal realized bodily demand when tissue and later motion diverge. All twenty futures complete; the maximum per-step purchase mismatch is $5.55\times10^{-17}$. The supply apparatus, sensor implementation, central budget coordination and computation remain incompletely priced. None of these branches is an independent replication of a new learned life.

\subsection*{Analysis and reproducibility}
Comparisons are descriptive and use the specified trajectory or complete-parent units. Opposite mappings, repeated phases, material sites and branches of one parent are not independent replicates. No population-level significance or general superiority is inferred from the two-parent mechanism tests. Missing outcomes are not imputed, later successes do not replace earlier failures, and realized costs are reported separately from behavioural criteria. A changed state can influence action without being the sole explanation of successful occupancy.

The seven main figures reuse their existing, editable presentation sources and verified simulation exports. No new simulation or illustration was generated for this editorial revision. Display scales, coordinate conventions, masks and paired comparisons are stated in the figure legends. Exact configurations, seeds and necessary initial conditions accompany the final code; figure data and original detailed protocols remain within this Article. These files retain the full precision and accounting needed beyond the rounded quantities in the main text.
